% F081: precise mathematical teaching diagram.
\begin{tikzpicture}[x=1cm,y=1cm,>=stealth,font=\small]
\draw[rounded corners=3pt,gray!40] (0,0) rectangle (14.2,8.6);
\begin{scope}[yshift=1cm]
\node at (7.1,7.0) {再用一次：三个零点逐层夹出两个、一个};
\draw[->,gray!70] (3,5.9)--(12.3,5.9);
\node[] at (1.55,5.9) {$f=0$};
\fill[blue!75!black] (3.5999999999999996,5.9) circle (2pt);\node[] at (3.5999999999999996,6.300000000000001) {$x_1$};
\fill[blue!75!black] (7.1,5.9) circle (2pt);\node[] at (7.1,6.300000000000001) {$x_2$};
\fill[blue!75!black] (10.6,5.9) circle (2pt);\node[] at (10.6,6.300000000000001) {$x_3$};
\draw[->,gray!70] (3,3.9)--(12.3,3.9);
\node[] at (1.55,3.9) {$f^{\prime}=0$};
\fill[blue!75!black] (5.07927405783631,3.9) circle (2pt);\node[] at (5.07927405783631,4.3) {$u$};
\fill[blue!75!black] (9.12072594216369,3.9) circle (2pt);\node[] at (9.12072594216369,4.3) {$v$};
\draw[->,gray!70] (3,1.9)--(12.3,1.9);
\node[] at (1.55,1.9) {$f^{\prime\prime}=0$};
\fill[blue!75!black] (7.1,1.9) circle (2pt);\node[] at (7.1,2.3) {$w$};
\draw[dashed,orange!75!black] (3.5999999999999996,5.75)--(5.07927405783631,4.55);
\draw[dashed,orange!75!black] (7.1,5.75)--(5.07927405783631,4.55);
\draw[dashed,orange!75!black] (7.1,5.75)--(9.12072594216369,4.55);
\draw[dashed,orange!75!black] (10.6,5.75)--(9.12072594216369,4.55);
\draw[dashed,orange!75!black] (5.07927405783631,3.75)--(7.1,2.55);
\draw[dashed,orange!75!black] (9.12072594216369,3.75)--(7.1,2.55);
\end{scope}
\node[text width=13cm,align=center] at (7.1,.9) {图取 $f=x^3-x$：$u=-1/\sqrt3,\ v=1/\sqrt3,\ w=0$。};
\end{tikzpicture}
